% !TeX root = ../papers/nc-p3.tex
\section{八余因子の楕円Bezout表示と四捩れ条件}
\label{sec:elliptic-bezout-torsion}
\begin{theorem}\label{thm:elliptic-bezout-torsion}
四辺有限な標準有限Frobenius窓の反対二次式行列 $N$ において $d=\det N$ が既約八次式、
余因子の線型包が八次元、その核曲線 $E$ が滑らかとする。
$L=\mathcal O_E(1)$ として $E$ は楕円正規四次曲線である。
前節の八次偏極 $L^2$ の四次鉛筆には、余因子への代入で恒等零となるただ一つの元がある。
それを $P=L\eta,Q=L\eta^{-1}$ の積の行列式で表す。
$\eta\in\operatorname{Pic}^0(E)$ は逆元まで一意で、次のどちらかである。
\begin{enumerate}
\item $\eta$ は非自明な二捩れで、鉛筆の分岐元である。
\item $\eta$ の位数は4である。この場合、入力基底変更で $N$ は
$P,Q$ の四切断空間の間の同型 $T$ による楕円Bezout行列
$\mathcal B(a,Ta)$ となる。行列が階数1となる三つの共通零点の因子では、
二つの切断の残余点は $\eta^2$ の平行移動だけ異なる。
\end{enumerate}
$H^0(\mathcal I_X^2(6))$ の次元や三重消滅八次式の一意性は仮定しない。
\end{theorem}
\begin{proof}[鉛筆の非分岐元の二つの核]
八個の余因子座標の空間だけを使えばよい。$\operatorname{adj}N$ は $D=V(d)$ で
核点の平方を与え、その像は $|L^2|$ の楕円正規八次曲線を張る。
前節の四つの対称行列式はこの像への代入でいずれも $d^3$ の定数倍である。
元の $L$ では非零であり、四つの元が一本の鉛筆を張るので消滅元はただ一つである。
全 $\eta$ の行列式からの写像は符号まで一点を分離する射影直線から射影直線への射であり、
この元も得られる。

$P\ne Q$ とし、余因子への代入で得る $H^0(P),H^0(Q)$ の対合を $C(x)$ とする。
一般の $x$ で $L,L$ の対合は可逆である。このとき $C$ は階数3である。
実際左核に独立な二切断があれば、その共通固定因子 $A$ の次数 $e$ は2以下である。
この二切断と $H^0(Q)$ の積写像の核は、基点を除いた鉛筆の評価完全列より
$H^0(QP^{-1}(A))$ となる。$e=0$ なら零で全八切断を張り、余因子汎関数が零になってしまう。
$e>0$ なら核次元は $e$ で、積は $A$ で消える八次切断を全て張る。
汎関数は長さ $e$ の $A$ 上の評価を経由し、$L,L$ 対合も階数 $e$ 以下になり矛盾する。

$D$ の一般点では $C$ は一点評価の積で階数1なので、$d^2$ はその随伴の全成分を割る。
商は一般階数1の非零な二次式行列である。左・右の核のどちらも定数ではない。
一つの固定切断と相手の全切断との積は非零であるが、八個の余因子座標が全空間を張るためである。
一意分解による階数1分解で、従って両方が入力 $x$ の一次式列 $a(x),b(x)$ となる。
一般の二切断は互いに素である。共通因子の次数 $e>0$ なら $P\ne Q$ より $e\le3$ で、
二つの積空間 $aH^0(Q)$ と $bH^0(P)$ の共通部分は
$ab$ 倍した $H^0(\mathcal O_E(A))$ の次元 $e$ である。
先と同様、汎関数は $A$ 上に支えられて可逆性に反する。
従って一般に二つの積空間の和は七次元となり、その零化直線が余因子汎関数である。
\end{proof}

\begin{proof}[対称なBezout行列]
$E\times E$ 上で $p_1^*\eta^{-1}p_2^*\eta$ に値を取り対角に単純極を許す核
$K(z,w)$ を取る。不変微分を固定し対角の留数を1とすれば一意に存在する。
実際、捩り前の直線束は非自明な次数零因子の外積で全コホモロジーが零であり、
対角の法線は自明である。
\[
 \mathcal B(a,b)(z,w)
   = a(z)b(w)K(z,w)+a(w)b(z)K(w,z)
\]
では対角の極が相殺され、$H^0(L)\otimes H^0(L)$ の対称元となる。
単純な四零点を持つ $a$ ではそれらでの $L$ の評価は基底である（$L/P$ は非自明な次数零）。
この基底の双対側で行列を評価すると、相異なる二点の成分は零であり、
各対角成分は固定した微分と自明化での $a'(z)b(z)$ である（全体の符号を除く）。
よって共通根のない一般対では可逆である。逆行列の定める八次汎関数は
この四点での評価の $1/(a'(z)b(z))$ 倍の和となる。
$aH^0(Q)$ は明らかに消え、$bH^0(P)$ での零は $h/a$ に不変微分を掛けた
楕円曲線上の留数の和が零なことによる。

この逆行列の射影公式は全ての互いに素な対で成立する。
実際、二つの切断に共通根があればその評価が核に入る。
共通根を印すと二つの平面束のファイバー積となり、像は既約で双次数 $(4,4)$ の因子である。
一方の一般切断の四根それぞれに対し、相手には相異なる一超平面を課すためである。
同じ双次数で零でない行列式の零因子はこの結果式に等しい。
上の逆行列と積の零化直線の等号は開集合上の正則写像の等号として延びる。
従って $N(x)$ は $\mathcal B(a(x),b(x))$ に射影的に等しい。
前者の十成分は二次式の基底で共通因子がなく、両者とも次数2なので全体定数しか違わない。
$a,b$ のいずれの線型写像にも核はない。さもなければ共通基点になる。
これで二つの四空間を同定して表示 $b=Ta$ を得る。
\end{proof}

\begin{proof}[残余点の入射と長さ十六]
$\delta\in E$ を $\eta^2$ に対応する点とし、加法記法を使う。
$N$ が $X$ 上で階数1なら、$a,b$ の共通零因子は $X$ 上の相対長さ3の因子となる。
非被約な場合を確認する。零因子 $D_a$ 上での $H^0(L)$ の評価は常に同型である。
$D_a$ 上で消えない補助切断 $b_0$ を取れば、この同定のもとで
$\mathcal B(a,b)$ は可逆な $\mathcal B(a,b_0)$ に $D_a$ 上の $b/b_0$ 倍を
出力で合成したものとなる。これは単純四根での対角公式から普遍の平坦な四点族上へ延びる。
単元小行列式を持つ階数1条件でその乗法の余核は底上階数3の局所自由加群となり、
$D_a\cap D_b$ が平坦長さ3である。滑らかな曲線の相対有限平坦部分スキームは有効相対因子であり
（局所の単生成をファイバーから持ち上げる）、残余は二点 $u,v=u-\delta$ の切断になる。

$u$ を動かす楕円曲線上の階数3束を
$A_u=H^0(P(-u)), A'_u=H^0(Q(-v))$ とする。
二つの線束 $P(-u),Q(-v)$ は各ファイバーで同型なので、相対的には
$A'=A\otimes J$ であり、$J$ の類は $-\delta$ である。
実際二つの対角因子（後者では $v=u-\delta$ のグラフ）を用いた普遍線束の比を
曲線因子の原点に制限すれば $\mathcal O([o]-[\delta])$ となる。
残余を取った共通三因子の二切断は単元スカラーを除いて等しい。
このスカラーを印す底を $\Sigma=\mathbb P(\mathcal O\oplus J)$ とする
（双対規約なら $J^{-1}$）。各 $u$ で二つの単射
$TA_u\hookrightarrow H^0(Q)$ と $A'_u\hookrightarrow H^0(Q)$ の差の
四行三列鉛筆を取る。正確には初めの写像は $A_u$ への $T$ の制限である。
両境界切断では単射なので、退化軌跡 $Z$ はスカラーが非零有限なところに入る。
$X$ はその射影化された核の入射に入る。入力 $a$ で復元できるので閉埋込みである。

$Z$ が有限のとき $X$ の長さは高々16である。一般の同型 $T$ での数から示す。
$\mathbb P(A)$ の共通三因子の二表示を二つの $\mathbb P^3$ に送る。
両射影の超平面類 $h_1,h_2$ は底からの次数零線束しか違わない。
両射影の次数は4で、一つの一般切断のどの零点を残すかの四選択に等しい。
従って対角類の一般移動との交点数は
$\int(h_1^3+h_1^2h_2+h_1h_2^2+h_2^3)=16$。
一般移動では有限横断的であり、各核直線からスカラーを一意に読む。

特殊な $T$ で $Z$ が有限でも核次元が跳ぶため、次の長さ比較が要る。
$\Sigma$ の四行三列写像を $F\to G$ とし、射影核上の二次式を表す複体
\[
 \wedge^2G^*\longrightarrow G^*\otimes F^*
     \longrightarrow\operatorname{Sym}^2 F^*
\]
を取る（鉛筆の直線束は $F,G$ に含める）。
一般階数3の開集合では三変数の基底と余分な零一本のKoszul複体の次数2なので完全である。
非完全軌跡の余次元は2以上で、階数 $6,6$ の二写像の完全性判定により
最後の有限長余核を除いて完全である。余核長は同じ束のEuler数だけで決まり、
一般横断の場合の16である。実際一般の各点では二つの単元行列を消すと、
残る一変数にかかる二局所方程式の商が余核になる。
階数0なら $A_u$ の射影平面全体が $N$ の階数1となり、整八次曲面が平面を含んでしまう。
従って各スカラー上で $X$ のファイバーは一点または入力の直線上にある。
直線切断長は非被約でも3以下なので、ファイバーでは二次式が全切断を張る。
相対中山の補題で最後の余核から $X$ 上の相対二次線束の押出しへ全射となり、
$20\le16$ という矛盾を得る。
\end{proof}

\begin{proof}[水平曲線と三捩れの排除]
従って $Z$ は曲線を含む。一つの $u$ の上の鉛筆は両端で単射なので垂直曲線はない。
水平曲線の正規化上ではスカラーが $J$ の消えない切断を与え、
底への有限写像のノルムで $J$ の位数が被覆次数を割る。
一般ファイバーで少なくとも一つの非零三次小行列式があり、その次数は3だから
$\delta$ の位数は2か3である。

位数3とすると、$J^3$ を自明化して水平曲線上のスカラーの三乗は非零定数となる。
各一般ファイバーで三根全てが現れる（正規化の射影次数は3の倍数であり、
また $J$ の全空間内の三乗根被覆は位数がちょうど3のため連結である）。
全ての最大小行列式は同じ三根を持つので、定数項と最高次項の四成分は互いに比例する。
二つの超平面のPl\"ucker座標であり $T A_u=A'_u$ が全 $u$ で成立する。
従って $T$ は評価の曲線を $u\mapsto u-\delta$ で運ぶ、平行移動の線型表示に等しい。
曲線は射影枠を含むので全体定数以外の自由度はない。

各 $u$ の三次因子の切断空間ではこの平行移動は射影位数3で、
三つの固有直線が全て相異なる。実際次数3の線束は $\delta$ の移動で不変で、
同型の三乗を1に規格化すると、位数3のエタール商へ次数1の線束として降りる。
三つの指標はその次数1線束の三つの零次捩りで、それぞれ切断が一次元である。
従ってスカラー付きの固有直線の入射は底 $u$ の上でエタール、特に被約である。
その三次零因子は $\delta$ 不変であり、一つの三点軌道全体である。
軌道を与えれば、残余 $u$ は全因子の線束 $P$ から一意に決まるため、
この入射は $E/\langle\delta\rangle$ の曲線そのものである。
入力超平面の次数は4である。固定点を含む軌道は一つで、残余点が指定点になるのは
三通りである。後者は $E$ 上の三倍写像の次数9を商の次数3で割ったものである。
よってこの曲線上の二次切断は八次元に過ぎず、入力三空間の十次元の二次式の
少なくとも二本が入射全体で消える。そこに入る $X$ の二次評価単射性に反する。
これで非分岐の場合は $\delta=2\eta$ の位数が2、すなわち $\eta$ の位数が4となる。
\end{proof}

\path{computations/verify_elliptic_bezout_torsion.py} は $y^2=x^3-x+1$ 上で
$(0,1),(0,-1)$ の切断を用い、二つの積空間の零化直線と逆行列から二次式表示を復元する。
整数同型の標本の良い還元では、二次成分の階数10、六次余因子の階数8である一方、
四次の二次小行列式は十九独立となる。有限辺の二十点を通る四次式は十五次元のため、
この標本は平方表示の次数だけを満たしても有限窓にならない。
